\documentclass[11pt]{article}
\usepackage[utf8]{inputenc}
\usepackage[T1]{fontenc}
\usepackage{geometry}
\usepackage{hyperref}
\usepackage{enumitem}
\geometry{letterpaper, margin=1in}
\title{\textbf{Working Draft Outline: Paper 1 \\ Conventional Massless Route}}
\author{Stephen Breighner}
\date{March 2026}
\begin{document}
\maketitle

\section*{Status}
This is a working outline for the conservative branch of the Bipolar Graviton Condensate project. The aim is to keep the strongest claims narrow: use a massless graviton baseline, treat the condensate as an effective description under extreme collapse, and defer the conversion mechanism to later work.

\section*{Core Scope}
\begin{itemize}[leftmargin=1.25em]
  \item Keep the long-range gravitational baseline close to the standard massless picture.
  \item Do not claim ordinary matter literally becomes gravitons in this paper.
  \item Present the condensate as an effective non-singular end-state ansatz.
  \item Use the toy potential and rotation-curve modifications as model outputs, not proofs of microscopic reality.
\end{itemize}

\section*{Proposed Flow}
\begin{enumerate}[leftmargin=1.4em]
  \item State the singularity problem and why an effective finite-radius core is worth modeling.
  \item Introduce the bipolar condensate ansatz as a phenomenological alternative.
  \item Define the effective potential and show the conditions for a stable radius.
  \item Compare the fringe velocity profile against the Schwarzschild baseline.
  \item Close with explicit limits: this is a consistency model, not yet a derivation.
\end{enumerate}

\section*{Required Clarifications}
\begin{itemize}[leftmargin=1.25em]
  \item Explain the support term as an effective short-distance regularization, not a derived interaction.
  \item Keep the language around ``dipoles'' and ``gravitons'' clearly marked as model assumptions.
  \item Distinguish singularity avoidance in the toy potential from a full GR result.
\end{itemize}

\section*{Near-Term Deliverable}
A clean first paper that argues: if a bipolar graviton condensate exists as an effective phase, then a finite-radius core and a modified fringe profile follow naturally from the chosen interaction ansatz.

\end{document}
